% ============================================================
% Nodes and Nations — A0 Portrait Academic Poster
% Compile with: pdflatex poster_a0.tex
% Requires: beamerposter, tikz, booktabs, colortbl, palatino
% ============================================================
\documentclass[final]{beamer}

% A0 Portrait: 841mm x 1189mm, scale=1.2 scales all fonts
\usepackage[
  orientation=portrait,
  size=a0,
  scale=1.2
]{beamerposter}

\usepackage{lmodern}
\usepackage{palatino}
\usepackage{graphicx}
\usepackage{booktabs}
\usepackage{colortbl}
\usepackage{array}
\usepackage{multirow}
\usepackage{amsmath,amssymb}
\usepackage{tikz}
\usetikzlibrary{arrows.meta, positioning, fit, backgrounds, decorations.pathmorphing}
\usepackage{tcolorbox}
\tcbuselibrary{skins,breakable}

% ── Colour palette ─────────────────────────────────────────
\definecolor{NavyBlue}   {RGB}{26,  42,  84}
\definecolor{SteelBlue}  {RGB}{52,  90, 147}
\definecolor{AccentRed}  {RGB}{192,  44,  26}
\definecolor{AccentGold} {RGB}{220, 160,  30}
\definecolor{LightGray}  {RGB}{245, 246, 248}
\definecolor{MidGray}    {RGB}{120, 130, 145}
\definecolor{BlockTitle} {RGB}{36,  67, 120}
\definecolor{BlockBody}  {RGB}{250, 251, 253}

% ── Beamer colour/font setup ────────────────────────────────
\usefonttheme{professionalfonts}
\setbeamercolor{background canvas}{bg=LightGray}
\setbeamercolor{block title}      {bg=BlockTitle, fg=white}
\setbeamercolor{block body}       {bg=BlockBody,  fg=black}
\setbeamerfont{block title}       {size=\large,series=\bfseries}
\setbeamerfont{block body}        {size=\normalsize}

% Remove navigation symbols
\setbeamertemplate{navigation symbols}{}

% ── Block template ──────────────────────────────────────────
\setbeamertemplate{block begin}{%
  \vskip 0.5ex
  \begin{beamercolorbox}[
      wd=\textwidth,
      ht=3ex,
      dp=1.5ex,
      leftskip=0.8em,
      rightskip=0.8em
    ]{block title}%
    \usebeamerfont{block title}\insertblocktitle
  \end{beamercolorbox}%
  \nointerlineskip
  \begin{beamercolorbox}[
      wd=\textwidth,
      sep=0.8em,
      leftskip=0.3em,
      rightskip=0.3em
    ]{block body}%
    \usebeamerfont{block body}%
}
\setbeamertemplate{block end}{%
  \end{beamercolorbox}%
  \vskip 1ex
}

% ── Headline ────────────────────────────────────────────────
\setbeamertemplate{headline}{%
  % Dark navy banner
  \begin{beamercolorbox}[wd=\paperwidth]{headline}
    \begin{tikzpicture}[remember picture]
      % Background fill
      \fill[NavyBlue] (0,0) rectangle (\paperwidth, 18cm);
      % Gold accent bar at bottom
      \fill[AccentGold] (0,0) rectangle (\paperwidth, 0.6cm);
      % Red accent bar
      \fill[AccentRed] (0,0.6cm) rectangle (\paperwidth, 1.1cm);

      % IIT Bhilai branding text (left)
      \node[anchor=west, text=white!80!NavyBlue,
            font=\large\itshape] at (2cm, 15cm)
        {Indian Institute of Technology Bhilai};
      \node[anchor=west, text=white!60!NavyBlue,
            font=\normalsize] at (2cm, 13.5cm)
        {Department of Mathematics and Computing};

      % Conference / course label (right)
      \node[anchor=east, text=AccentGold,
            font=\large\bfseries] at (\paperwidth-2cm, 15cm)
        {Network Science · CS6XX · 2026};
      \node[anchor=east, text=white!70!NavyBlue,
            font=\normalsize] at (\paperwidth-2cm, 13.5cm)
        {Nodes and Nations Project};

      % Main title
      \node[anchor=north, text=white, align=center,
            font=\fontsize{52}{60}\selectfont\bfseries] at (0.5*\paperwidth, 12cm)
        {Nodes and Nations};
      \node[anchor=north, text=white!85!NavyBlue, align=center,
            font=\fontsize{32}{40}\selectfont] at (0.5*\paperwidth, 7.8cm)
        {A Complex Network Study of Global Migration (1990–2025)};

      % Authors
      \node[anchor=north, text=AccentGold, align=center,
            font=\fontsize{28}{36}\selectfont\bfseries] at (0.5*\paperwidth, 5.6cm)
        {Aditya Yadav \quad Dipanjan Mondal \quad Jaswinder Singh \quad Praveen Kumar};

      % Horizontal rule
      \draw[white!40!NavyBlue, line width=1pt]
        (2cm, 3.2cm) -- (\paperwidth-2cm, 3.2cm);
      \node[anchor=north, text=white!70!NavyBlue, align=center,
            font=\fontsize{22}{28}\selectfont\itshape] at (0.5*\paperwidth, 3.0cm)
        {UN DESA Bilateral Migrant Stock • World Bank WDI •
         UCDP/PRIO • UNDP HDR • ND-GAIN};
    \end{tikzpicture}%
  \end{beamercolorbox}%
}

% ── Helpers ─────────────────────────────────────────────────
\newcommand{\colsep}{2cm}   % gap between columns

% Highlighted box (for key findings)
\newcommand{\keybox}[2]{%
  \begin{tcolorbox}[
      enhanced,
      colback=NavyBlue!8,
      colframe=AccentRed,
      leftrule=4pt, rightrule=0pt, toprule=0pt, bottomrule=0pt,
      boxsep=4pt, left=6pt, right=6pt,
      title={\bfseries\color{AccentRed} #1},
      fonttitle=\normalsize,
      coltitle=AccentRed,
      attach boxed title to top left={yshift=-2mm,xshift=6pt},
      boxed title style={colback=white,colframe=AccentRed,
                         leftrule=0pt,rightrule=0pt,
                         toprule=0pt,bottomrule=0pt}
    ]
    #2
  \end{tcolorbox}%
}

% ─────────────────────────────────────────────────────────────
\begin{document}
\begin{frame}[t,fragile]

% Outer padding
\hspace{1.5cm}
\begin{columns}[T]

% ===========================================================
% COLUMN 1  (left)
% ===========================================================
\begin{column}{0.30\textwidth}

  % --- Abstract / Motivation --------------------------------
  \begin{block}{1 · Motivation \& Research Questions}
    Migration is not random. It expresses itself as a
    \textbf{directed, weighted network} in which countries are nodes
    and mass human movement is an edge whose weight encodes
    bilateral migrant stock.

    \vspace{0.6em}
    \textbf{Key Research Questions:}
    \begin{enumerate}
      \item Which nations act as \emph{global hubs}
            (high Betweenness) vs.\ \emph{terminal attractors}
            (high PageRank)?
      \item Do geopolitical communities become more or
            less compartmentalised over 35 years?
      \item Can standard macroeconomic variables (GDP,
            population, conflict) explain observed edge weights?
    \end{enumerate}
  \end{block}

  % --- Data -------------------------------------------------
  \begin{block}{2 · Dataset \& Pipeline}
    \textbf{Primary Migration Data}\\
    UN DESA International Migrant Stock (2024 release):
    $\approx$\,286 origin $\times$ 286 destination countries;
    5-year intervals 1990–2020 extended to 2025 via
    log-linear extrapolation.

    \vspace{0.7em}
    \textbf{Supplementary Factor Datasets}
    \begin{center}
    \small
    \renewcommand{\arraystretch}{1.3}
    \begin{tabular}{>{\bfseries}l l}
      \toprule
      Source & Variable \\
      \midrule
      World Bank WDI & GDP per Capita, Population,
                        Unemployment \\
      UCDP/PRIO ACD  & Conflict Intensity \\
      UNDP HDR       & Education Index (MYS) \\
      ND-GAIN (2026) & Climate Vulnerability Score \\
      \bottomrule
    \end{tabular}
    \end{center}

    \vspace{0.7em}
    \textbf{Processing Pipeline}
    \vspace{0.5em}

    % Simple horizontal pipeline with TikZ
    \begin{center}
    \begin{tikzpicture}[
        box/.style={draw=SteelBlue, fill=SteelBlue!12,
                    rounded corners=4pt, minimum width=2.4cm,
                    minimum height=1.1cm, align=center,
                    font=\small\bfseries},
        arr/.style={-Stealth, SteelBlue, thick},
        node distance=0.35cm
      ]
      \node[box] (a) {01\\Ingest \&\\Clean};
      \node[box, right=of a] (b) {02\\Network\\Build};
      \node[box, right=of b] (c) {03\\Community\\Detect};
      \node[box, right=of c] (d) {04\\OLS\\Regression};
      \node[box, right=of d] (e) {05\\Power BI\\Export};
      \draw[arr] (a)--(b); \draw[arr] (b)--(c);
      \draw[arr] (c)--(d); \draw[arr] (d)--(e);
    \end{tikzpicture}
    \end{center}

    \vspace{0.5em}
    \underline{Outputs:}
    \texttt{migration\_long.csv} (54,521 rows),
    \texttt{factors\_panel.csv} (2,130 rows),
    centrality and community CSVs, plus a Power BI
    star-schema manifest.
  \end{block}

  % --- Phase 2 centrality -----------------------------------
  \begin{block}{3 · Network Topology \& Centrality}

    We construct \textbf{8 directed, weighted snapshots}
    (NetworkX). Edge weight = bilateral migrant stock.

    \vspace{0.5em}
    \textbf{Centrality Measures Applied:}
    \begin{itemize}
      \item \textbf{In/Out-Degree Centrality} — raw
            connection count normalised by $N-1$.
      \item \textbf{Betweenness Centrality} — fraction of
            shortest paths passing through node $v$:
            $C_B(v)=\sum_{s\neq v\neq t}
             \frac{\sigma_{st}(v)}{\sigma_{st}}$
      \item \textbf{Weight-aware PageRank} — iterative
            authority score proportional to incoming
            \emph{migrant volume}.
    \end{itemize}

    \vspace{0.5em}
    \keybox{Key Finding}{%
      High-betweenness \emph{transit hubs}
      (e.g.\ Russia, UAE, Jordan) are structurally distinct
      from high-PageRank \emph{economic attractors}
      (e.g.\ USA, Germany, Australia). Both sets expanded
      over time, confirming increasing network concentration.
    }

  \end{block}

\end{column}

% ===========================================================
% COLUMN 2  (centre)
% ===========================================================
\hspace{\colsep}
\begin{column}{0.30\textwidth}

  % --- Community detection -----------------------------------
  \begin{block}{4 · Community Detection}

    Three algorithms applied to each snapshot graph:

    \vspace{0.4em}
    \begin{center}
    \begin{tikzpicture}[
        algo/.style={draw=NavyBlue, fill=NavyBlue!10,
                     rounded corners=4pt, minimum width=3.5cm,
                     minimum height=1cm, align=center,
                     font=\small\bfseries},
        arr/.style={-Stealth, NavyBlue, thick}
      ]
      \node[algo] (lo) {Louvain};
      \node[algo, right=1.2cm of lo] (le) {Leiden};
      \node[algo, right=1.2cm of le] (gn) {Girvan-Newman};
      \node[draw=AccentGold, fill=AccentGold!20, rounded corners=4pt,
            minimum width=10cm, minimum height=1cm,
            align=center, below=1cm of le,
            font=\small\bfseries] (out)
        {Communities + Boundary Nodes + Temporal Drift (Jaccard)};
      \draw[arr] (lo.south) -- (out.north west);
      \draw[arr] (le.south) -- (out.north);
      \draw[arr] (gn.south) -- (out.north east);
    \end{tikzpicture}
    \end{center}

    \vspace{0.8em}
    \textbf{Louvain Modularity $Q$ — Temporal Decline}

    \vspace{0.5em}
    \begin{center}
    % Handcrafted bar chart using TikZ (real data)
    \begin{tikzpicture}[xscale=1.05, yscale=11]
      % Axes
      \draw[->] (0,0.44) -- (0,0.585) node[above, font=\small]{$Q$};
      \draw[->] (0,0.44) -- (9.2,0.44) node[right, font=\small]{Year};
      % Grid lines
      \foreach \y/\label in {0.45/0.45,0.48/0.48,0.51/0.51,0.54/0.54,0.57/0.57}{%
        \draw[gray!40, dashed] (0,\y) -- (8.8,\y);
        \node[left, font=\tiny] at (0,\y) {\label};
      }
      % Bars (real Q values from modularity_scores.csv, Louvain only)
      % Year -> x offset
      \foreach \x/\yr/\q in {
        0.5/1990/0.5595,
        1.6/1995/0.5489,
        2.7/2000/0.5321,
        3.8/2005/0.5106,
        4.9/2010/0.5013,
        6.0/2015/0.4930,
        7.1/2020/0.4737,
        8.2/2025/0.4748%
      }{%
        \fill[SteelBlue!80] (\x,0.44) rectangle (\x+0.8,\q);
        \node[above, font=\tiny\bfseries, color=NavyBlue]
          at (\x+0.4, \q) {\q};
        \node[below, font=\tiny] at (\x+0.4, 0.44) {\yr};
      }
      % Trend line (approximate)
      \draw[AccentRed, ultra thick, dashed]
        (0.9,0.5595) -- (9.0, 0.4748);
    \end{tikzpicture}
    \end{center}

    \keybox{Key Finding}{%
      Louvain modularity fell \textbf{0.560 → 0.475}
      (1990–2025): the global migration graph is becoming
      \emph{less compartmentalised} — regional clusters are
      dissolving into an integrated, hub-centric topology.
    }

    \vspace{0.8em}
    \textbf{Temporal Structural Drift (Jaccard Similarity)}

    \vspace{0.5em}
    \begin{center}
    \small
    \renewcommand{\arraystretch}{1.2}
    \begin{tabular}{cc c}
      \toprule
      \textbf{Period} & \textbf{Jaccard} & \textbf{Interpretation} \\
      \midrule
      1990–1995 & 0.702 & Moderate stability \\
      1995–2000 & 0.610 & Post-Soviet disruption \\
      2000–2005 & 0.854 & High continuity \\
      2005–2010 & 0.477 & \cellcolor{AccentRed!20}
                           \textbf{Largest shift} \\
      2010–2015 & 0.706 & Recovery \\
      2015–2020 & 0.536 & Refugee crisis impact \\
      2020–2025 & 0.635 & Pandemic rebound \\
      \bottomrule
    \end{tabular}
    \end{center}

    \vspace{0.5em}
    Low Jaccard $\Rightarrow$ many \textbf{boundary nodes}
    switched cluster affiliation, signalling geopolitical
    re-alignment.

  \end{block}

  % --- Algorithm Comparison ---------------------------------
  \begin{block}{5 · Algorithm Comparison (2020 Snapshot)}
    \begin{center}
    \small
    \renewcommand{\arraystretch}{1.2}
    \begin{tabular}{lccc}
      \toprule
      \textbf{Algorithm} & \textbf{$Q$} &
      \textbf{Communities} & \textbf{NMI vs.\ Louvain} \\
      \midrule
      Louvain       & 0.474 & 6  & — \\
      Leiden        & 0.092 & 8  & 0.828 \\
      Girvan-Newman & 0.021 & 64 & 0.345 \\
      \bottomrule
    \end{tabular}
    \end{center}

    \vspace{0.4em}
    High NMI (0.828) between Louvain and Leiden confirms
    \emph{robust} cluster structure. Girvan-Newman yields
    fine-grained sub-communities useful for boundary-node
    identification.
  \end{block}

\end{column}

% ===========================================================
% COLUMN 3  (right)
% ===========================================================
\hspace{\colsep}
\begin{column}{0.30\textwidth}

  % --- Regression -------------------------------------------
  \begin{block}{6 · Regression Analysis}

    \textbf{Gravity-inspired OLS model} fitted on
    panel data merging centrality metrics with
    the factors panel (n~=~1,448 country-year observations).

    \vspace{0.6em}
    \textbf{Baseline (3-factor) specification:}
    \[
      \ln W_{it} = \alpha + \beta_1\ln\text{Pop}_{it}
                   + \beta_2\ln\text{GDP}_{it}
                   + \beta_3\,\text{Conflict}_{it}
                   + \varepsilon_{it}
    \]

    \textbf{Full (6-factor) model adds:}
    Unemployment, Education Index (UNDP MYS),
    Climate Vulnerability (ND-GAIN).

    \vspace{0.8em}
    \textbf{Model Comparison}

    \vspace{0.4em}
    \begin{center}
    \small
    \renewcommand{\arraystretch}{1.25}
    \begin{tabular}{lcccc}
      \toprule
      \textbf{Model} & \textbf{$n$} &
      \textbf{$R^2$} & \textbf{Adj.~$R^2$} & \textbf{AIC} \\
      \midrule
      Baseline Inflows  & 1\,243 & 0.636 & 0.635 & 4156.7 \\
      Full Inflows      &   974  & 0.593 & 0.591 & 3186.3 \\
      Baseline Outflows & 1\,243 & 0.692 & 0.691 & 3961.0 \\
      \rowcolor{SteelBlue!18}
      \textbf{Full Outflows} & \textbf{974} &
      \textbf{0.708} & \textbf{0.706} & \textbf{2657.1} \\
      \bottomrule
    \end{tabular}
    \end{center}

    \vspace{0.6em}
    \textbf{VIF (full model, pooled):}

    \begin{center}
    \small
    \renewcommand{\arraystretch}{1.2}
    \begin{tabular}{lcc}
      \toprule
      \textbf{Predictor} & \textbf{VIF} & \textbf{Flag} \\
      \midrule
      $\ln$Population     & 1.32 & \textcolor{teal}{\checkmark} \\
      $\ln$GDP per Capita & 3.13 & \textcolor{teal}{\checkmark} \\
      Conflict Intensity  & 1.16 & \textcolor{teal}{\checkmark} \\
      Unemployment        & 1.16 & \textcolor{teal}{\checkmark} \\
      Education Index     & 3.23 & \textcolor{teal}{\checkmark} \\
      Climate Vulnerability & 4.32 & \textcolor{teal}{\checkmark} \\
      \bottomrule
    \end{tabular}
    \end{center}
    \vspace{0.3em}
    All VIF $<$ 10; no harmful multicollinearity detected.

    \keybox{Key Finding}{%
      GDP per capita and population alone account for
      \textbf{$>$63\,\%} of migration-flow variance.
      Adding climate and education raises outflow explanation
      to \textbf{70.6\,\%} (Adj.~$R^2=0.706$).
    }

  \end{block}

  % --- Conclusions ------------------------------------------
  \begin{block}{7 · Conclusions}
    \begin{enumerate}
      \item \textbf{Network Integration}: Louvain $Q$ declined
            from 0.560 (1990) to 0.475 (2025) — the world
            migration graph is structurally \emph{converging}.
      \item \textbf{Hub Consolidation}: Betweenness and
            PageRank leaders intensified their roles over time;
            routing is increasingly hub-mediated.
      \item \textbf{Predictable Drivers}: Standard gravity
            variables (GDP, population, conflict) explain
            $\sim$70\,\% of outflow variance, confirming
            macro-determinism in global mobility.
      \item \textbf{Boundary-Node Dynamics}: Post-2005
            geopolitical re-alignment produced the largest
            structural break (Jaccard = 0.477), corresponding
            to the EU expansion and GCC labour boom.
    \end{enumerate}
  \end{block}

  % --- Future Work ------------------------------------------
  \begin{block}{8 · Future Work}
    \begin{itemize}
      \item Integrate time-varying \textbf{visa-openness index}
            for improved bilateral modelling.
      \item Apply \textbf{Temporal Graph Networks (TGN)} for
            link prediction of future migration corridors.
      \item Extend Power BI dashboard with \textbf{drill-through
            node profiles} and animated temporal sliders.
      \item Explore \textbf{causal inference} (DiD, IV) to
            isolate conflict shocks from trend.
    \end{itemize}
  \end{block}

  % --- References -------------------------------------------
  \begin{block}{References}
    \small
    \begin{enumerate}
      \item UN DESA (2024). \textit{International Migrant Stock}.
      \item World Bank. \textit{World Development Indicators}.
      \item UCDP/PRIO. \textit{Armed Conflict Dataset v25.1}.
      \item UNDP (2024). \textit{Human Development Report}.
      \item ND-GAIN (2026). \textit{Country Index}.
      \item Blondel et al.\ (2008). Fast unfolding of communities
            in large networks. \textit{J.\ Stat.\ Mech.}
      \item Page et al.\ (1999). \textit{The PageRank Citation
            Ranking}. Stanford TR.
    \end{enumerate}
  \end{block}

\end{column}

\end{columns}

% ── Footer ──────────────────────────────────────────────────
\vspace{1cm}
\begin{beamercolorbox}[wd=\paperwidth, ht=2cm, dp=0.5cm]{headline}
  \begin{tikzpicture}
    \fill[NavyBlue] (0,0) rectangle (\paperwidth, 2.5cm);
    \fill[AccentGold] (0,2.5cm) rectangle (\paperwidth, 2.8cm);
    \node[text=white!80!NavyBlue, font=\normalsize,
          anchor=west] at (2cm, 1.25cm)
      {\textbf{IIT Bhilai} · Network Science · 2026 \quad
       Codebase: \texttt{github.com/[repo]/nodes-and-nations}};
    \node[text=AccentGold, font=\normalsize\bfseries,
          anchor=east] at (\paperwidth-2cm, 1.25cm)
      {Contact: \texttt{\{aditya, dipanjan, jaswinder, praveen\}@iitbhilai.ac.in}};
  \end{tikzpicture}
\end{beamercolorbox}

\end{frame}
\end{document}
